% ---- score2 standard preamble ----
\documentclass[10pt]{extarticle}
\usepackage[a4paper,margin=1.5cm]{geometry}
\usepackage[T1]{fontenc}
\usepackage{amsmath,amssymb}
\usepackage{graphicx}
\usepackage[export]{adjustbox}
\usepackage{multicol}
\usepackage{enumitem}
\usepackage{booktabs,array,tabularx}
\setlength{\columnsep}{0.7cm}
\setlength{\parindent}{0pt}
\setlength{\parskip}{3pt}
\pagestyle{empty}
\makeatletter
\newcommand{\qfigdef}[4]{\@namedef{qf@file@#1}{#2}\@namedef{qf@w@#1}{#3}\@namedef{qf@wide@#1}{#4}}
\newcommand{\qf@check}[1]{\@ifundefined{qf@file@#1}{\PackageError{score2}{Unknown figure id #1}{}}{}}
\newcommand{\qf@img}[4][]{\includegraphics[width=#2\dimexpr\@nameuse{qf@w@#4}pt\relax,max width=#3,max height=0.45\textheight,#1]{\@nameuse{qf@file@#4}}}
\newif\ifqf@full
\newif\ifqf@boxed
\newsavebox{\qf@box}
\newcommand{\qf@open}{\global\qf@fulltrue\par\global\qf@boxedtrue
  \setbox\qf@box\vbox\bgroup\hsize\textwidth\linewidth\textwidth}
\newcommand{\qf@emit}[1]{\par\egroup\global\qf@boxedfalse
  \ifdim\dimexpr\ht\qf@box+\dp\qf@box\relax<#1\textheight\noindent\usebox{\qf@box}\par\else\unvbox\qf@box\fi}
\newcommand{\qfullsplit}{\ifqf@boxed\qf@emit{0.85}\fi}
\newcommand{\qf@close}{\ifqf@boxed\qf@emit{0.45}\fi\global\qf@fullfalse}
\newcommand{\qfullbegin}{\end{multicols}\qf@open}
\newcommand{\qfullend}{\qf@close\begin{multicols}{2}}
\newcommand{\qfullstart}{\qf@open}
\newcommand{\qfulldone}{\qf@close}
\newcommand{\qfig}[2][1]{\qf@check{#2}%
  \ifqf@full\qf@blk{#1}{#2}\else
  \ifnum\@nameuse{qf@wide@#2}=1 \expandafter\@firstoftwo\else\expandafter\@secondoftwo\fi
  {\end{multicols}\begin{center}\qf@img{#1}{\textwidth}{#2}\end{center}\begin{multicols}{2}}%
  {\qf@blk{#1}{#2}}\fi}
\newcommand{\qf@blk}[2]{\par\smallskip{\centering\qf@img{#1}{\linewidth}{#2}\par}\smallskip}
\newcommand{\qfigi}[2][1]{\qf@check{#2}\qf@img[valign=c]{#1}{\linewidth}{#2}}
\makeatother
\newcommand{\papersec}[1]{\par\bigskip{\centering\Large\bfseries PAPER-#1\par}\medskip}
\newcommand{\partsec}[2]{\par\medskip\noindent{\large\bfseries PART-#1: #2}\par\smallskip}
\newcommand{\sect}[1]{\par\smallskip\noindent\textbf{#1}\par}
\newcommand{\qn}[1]{\par\smallskip\noindent\textbf{#1.}~}
\newcommand{\paraq}[2]{\par\smallskip\noindent\textbf{Paragraph for Questions #1 and #2}\par}
\newenvironment{qoptions}{\begin{itemize}[leftmargin=1.8em,itemsep=1pt,topsep=2pt,label={}]}{\end{itemize}}
\newcommand{\opt}[2]{\item[(#1)] #2}
% figure registry (generated by score2_tools.py finalize)
\qfigdef{p1_phy_s1_q02_1}{Score-II - Test 12 (04 Feb 2024)-img/p1_phy_s1_q02_1.png}{89.7}{0}
\qfigdef{p1_phy_s2_q05_A}{Score-II - Test 12 (04 Feb 2024)-img/p1_phy_s2_q05_A.png}{263.5}{1}
\qfigdef{p1_phy_s2_q06_1}{Score-II - Test 12 (04 Feb 2024)-img/p1_phy_s2_q06_1.png}{145.8}{0}
\qfigdef{p1_phy_s2_q06_2}{Score-II - Test 12 (04 Feb 2024)-img/p1_phy_s2_q06_2.png}{134.2}{0}
\qfigdef{p1_phy_s2_q07_1}{Score-II - Test 12 (04 Feb 2024)-img/p1_phy_s2_q07_1.png}{134.9}{0}
\qfigdef{p1_phy_s3_q11_1}{Score-II - Test 12 (04 Feb 2024)-img/p1_phy_s3_q11_1.png}{89.7}{0}
\qfigdef{p1_phy_s4_q02_1}{Score-II - Test 12 (04 Feb 2024)-img/p1_phy_s4_q02_1.png}{158.9}{0}
\qfigdef{p1_phy_s4_q06_1}{Score-II - Test 12 (04 Feb 2024)-img/p1_phy_s4_q06_1.png}{76.5}{0}
\qfigdef{p1_chem_s1_q01_1}{Score-II - Test 12 (04 Feb 2024)-img/p1_chem_s1_q01_1.png}{195.6}{0}
\qfigdef{p1_chem_s2_q06_1}{Score-II - Test 12 (04 Feb 2024)-img/p1_chem_s2_q06_1.png}{174.6}{0}
\qfigdef{p1_chem_s2_q06_2}{Score-II - Test 12 (04 Feb 2024)-img/p1_chem_s2_q06_2.png}{132.2}{0}
\qfigdef{p1_chem_s2_q06_3}{Score-II - Test 12 (04 Feb 2024)-img/p1_chem_s2_q06_3.png}{132.7}{0}
\qfigdef{p1_chem_s2_q07_A}{Score-II - Test 12 (04 Feb 2024)-img/p1_chem_s2_q07_A.png}{128.0}{0}
\qfigdef{p1_chem_s2_q07_B}{Score-II - Test 12 (04 Feb 2024)-img/p1_chem_s2_q07_B.png}{133.0}{0}
\qfigdef{p1_chem_s2_q07_C}{Score-II - Test 12 (04 Feb 2024)-img/p1_chem_s2_q07_C.png}{133.4}{0}
\qfigdef{p1_chem_s2_q07_D}{Score-II - Test 12 (04 Feb 2024)-img/p1_chem_s2_q07_D.png}{134.7}{0}
\qfigdef{p1_chem_s3_q10_4}{Score-II - Test 12 (04 Feb 2024)-img/p1_chem_s3_q10_4.png}{105.8}{0}
\qfigdef{p1_chem_s4_q03_1}{Score-II - Test 12 (04 Feb 2024)-img/p1_chem_s4_q03_1.png}{127.6}{0}
\qfigdef{p1_chem_s4_q04_1}{Score-II - Test 12 (04 Feb 2024)-img/p1_chem_s4_q04_1.png}{143.7}{0}
\qfigdef{p1_chem_s4_q05_1}{Score-II - Test 12 (04 Feb 2024)-img/p1_chem_s4_q05_1.png}{153.0}{0}
\qfigdef{p1_chem_s4_q06_1}{Score-II - Test 12 (04 Feb 2024)-img/p1_chem_s4_q06_1.png}{164.3}{0}
\qfigdef{p2_phy_s1_q01_1}{Score-II - Test 12 (04 Feb 2024)-img/p2_phy_s1_q01_1.png}{173.4}{0}
\qfigdef{p2_phy_s2_q05_1}{Score-II - Test 12 (04 Feb 2024)-img/p2_phy_s2_q05_1.png}{156.2}{0}
\qfigdef{p2_phy_s2_q06_1}{Score-II - Test 12 (04 Feb 2024)-img/p2_phy_s2_q06_1.png}{169.2}{0}
\qfigdef{p2_phy_s3_para1_2_1}{Score-II - Test 12 (04 Feb 2024)-img/p2_phy_s3_para1_2_1.png}{93.0}{0}
\qfigdef{p2_phy_s3_para3_4_1}{Score-II - Test 12 (04 Feb 2024)-img/p2_phy_s3_para3_4_1.png}{245.2}{0}
\qfigdef{p2_phy_s4_q05_1}{Score-II - Test 12 (04 Feb 2024)-img/p2_phy_s4_q05_1.png}{157.2}{0}
\qfigdef{p2_phy_s4_q06_1}{Score-II - Test 12 (04 Feb 2024)-img/p2_phy_s4_q06_1.png}{105.6}{0}
\qfigdef{p2_phy_s4_q08_1}{Score-II - Test 12 (04 Feb 2024)-img/p2_phy_s4_q08_1.png}{148.0}{0}
\qfigdef{p2_phy_s4_q10_1}{Score-II - Test 12 (04 Feb 2024)-img/p2_phy_s4_q10_1.png}{263.7}{1}
\qfigdef{p2_chem_s1_q01_A}{Score-II - Test 12 (04 Feb 2024)-img/p2_chem_s1_q01_A.png}{121.8}{0}
\qfigdef{p2_chem_s1_q01_B}{Score-II - Test 12 (04 Feb 2024)-img/p2_chem_s1_q01_B.png}{166.5}{0}
\qfigdef{p2_chem_s1_q01_C}{Score-II - Test 12 (04 Feb 2024)-img/p2_chem_s1_q01_C.png}{165.6}{0}
\qfigdef{p2_chem_s1_q01_D}{Score-II - Test 12 (04 Feb 2024)-img/p2_chem_s1_q01_D.png}{161.5}{0}
\qfigdef{p2_chem_s1_q03_1}{Score-II - Test 12 (04 Feb 2024)-img/p2_chem_s1_q03_1.png}{103.8}{0}
\qfigdef{p2_chem_s1_q03_2}{Score-II - Test 12 (04 Feb 2024)-img/p2_chem_s1_q03_2.png}{105.6}{0}
\qfigdef{p2_chem_s1_q03_3}{Score-II - Test 12 (04 Feb 2024)-img/p2_chem_s1_q03_3.png}{126.9}{0}
\qfigdef{p2_chem_s1_q03_4}{Score-II - Test 12 (04 Feb 2024)-img/p2_chem_s1_q03_4.png}{119.0}{0}
\qfigdef{p2_chem_s1_q04_A}{Score-II - Test 12 (04 Feb 2024)-img/p2_chem_s1_q04_A.png}{166.2}{0}
\qfigdef{p2_chem_s1_q04_B}{Score-II - Test 12 (04 Feb 2024)-img/p2_chem_s1_q04_B.png}{154.6}{0}
\qfigdef{p2_chem_s1_q04_C}{Score-II - Test 12 (04 Feb 2024)-img/p2_chem_s1_q04_C.png}{159.8}{0}
\qfigdef{p2_chem_s1_q04_D}{Score-II - Test 12 (04 Feb 2024)-img/p2_chem_s1_q04_D.png}{144.8}{0}
\qfigdef{p2_chem_s2_q07_1}{Score-II - Test 12 (04 Feb 2024)-img/p2_chem_s2_q07_1.png}{188.9}{0}
\qfigdef{p2_chem_s2_q07_3}{Score-II - Test 12 (04 Feb 2024)-img/p2_chem_s2_q07_3.png}{81.5}{0}
\qfigdef{p2_chem_s3_para1_2_1}{Score-II - Test 12 (04 Feb 2024)-img/p2_chem_s3_para1_2_1.png}{157.0}{0}
\qfigdef{p2_math_s3_para3_4_1}{Score-II - Test 12 (04 Feb 2024)-img/p2_math_s3_para3_4_1.png}{87.5}{0}
\begin{document}
\begin{center}{\Large\bfseries JEE (Main+Advanced) Score-II --- Test 12 (04 Feb 2024)}\end{center}
\begin{multicols}{2}
\papersec{1}
\partsec{1}{PHYSICS}
\sect{Section-I (i): Single Correct}

\qn{1} A bar of mas $M=2$ kg and length $L=0.20$ m is lying on a horizontal frictionless surface. One end of the bar is pivoted at a point about which it is free to rotate. A small mass $m=0.10$ kg is moving on the same horizontal surface with $5.00\ ms^{-1}$ speed on a path perpendicular to the bar. It hits the bar at a distance $\dfrac L2$ from the pivoted end and returns back on the same path with speed v. After this elastic collision, the bar rotates with an angular velocity $\omega$. Which of the following statement is correct?
\begin{qoptions}
\opt{A}{$\omega=6.98$ rad $s^{-1}$ and $v=4.30$ m $s^{-1}$}
\opt{B}{$\omega=3.61$ rad $s^{-1}$ and $v=4.63$ m $s^{-1}$}
\opt{C}{$\omega=3.75$ rad $s^{-1}$ and $v=10.0$ m $s^{-1}$}
\opt{D}{$\omega=6.80$ rad $s^{-1}$ and $v=4.10$ m $s^{-1}$}
\end{qoptions}

\qn{2} A container has a base of 50 cm $\times$ 5 cm and height 50 cm, as shown in the figure. It has two parallel electrically conducting walls each of area 50 cm $\times$ 50 cm. The remaining walls of the container are thin and non-conducting. The container is being filled with a liquid of dielectric constant 3 at a uniform rate of 250 $cm^3\,s^{-1}$. What is the value of the capacitance of the container after 20 seconds?
[Given: Permittivity of free space $\epsilon_0=9\times10^{-12}\ C^2N^{-1}m^{-2}$, the effects of the non-conducting walls on the capacitance are negligible]
\qfig{p1_phy_s1_q02_1}
\begin{qoptions}
\opt{A}{27 pF}
\opt{B}{63 pF}
\opt{C}{81 pF}
\opt{D}{135 pF}
\end{qoptions}

\qn{3} One mole of an ideal gas expands adiabatically from an initial state $(T_A, V_0)$ to final state $(T_f, 5V_0)$. Another mole of the same gas expands isothermally from a different initial state $(T_B, V_0)$ to the same final state $(T_f, 5V_0)$. The ratio of the specific heats at constant pressure and constant volume of this ideal gas is $\gamma$. What is the ratio $T_A/T_B$?
\begin{qoptions}
\opt{A}{$5^{\gamma-1}$}
\opt{B}{$5^{1-\gamma}$}
\opt{C}{$5^{\gamma}$}
\opt{D}{$5^{1+\gamma}$}
\end{qoptions}

\qn{4} Two satellites P and Q are moving in different circular orbits around the Earth (radius R). The heights of P and Q from the Earth surface are $h_P$ and $h_Q$, respectively, where $h_P=\dfrac{R}{2}$. The accelerations of P and Q due to Earth's gravity are $g_P$ and $g_Q$, respectively. If $\dfrac{g_P}{g_Q}=\dfrac{49}{25}$, what is the value of $h_Q$?
\begin{qoptions}
\opt{A}{$\dfrac{3R}{5}$}
\opt{B}{$\dfrac{R}{6}$}
\opt{C}{$\dfrac{21R}{10}$}
\opt{D}{$\dfrac{11R}{10}$}
\end{qoptions}

\qfullbegin
\sect{Section-I (ii): One or More Correct}

\qn{5} A slide with a frictionless curved surface, which becomes horizontal at its lower end, is fixed on the terrace of a building of height $2h$ from the ground. A spherical ball of mass m is released on the slide from rest at a height h from the top of the terrace. The ball leaves the slide with a velocity $\vec u_0=u_0\hat x$ and falls on the ground at a distance d from the building making an angle $\theta$ with the horizontal. It bounces off with a velocity $\vec v$ and reaches a maximum height $h_1$. The acceleration due to gravity is g and the coefficient of restitution of the ground is $\dfrac{1}{\sqrt2}$. Which of the following statement(s) is(are) correct?
\qfig{p1_phy_s2_q05_A}
\qfullsplit
\begin{qoptions}
\opt{A}{$\vec u_0=\sqrt{2gh}\hat x$}
\opt{B}{$\vec v=\sqrt{2gh}(\hat x+\hat z)$}
\opt{C}{$\theta=60^\circ$}
\opt{D}{$\dfrac{d}{h_1}=2\sqrt3$}
\end{qoptions}
\qfullend

\qn{6} A plane polarized blue light ray is incident on a prism such that there is no reflection from the surface of the prism. The angle of deviation of the emergent ray is $\delta=60^\circ$ (see Figure-1). The angle of minimum deviation for red light from the same prism is $\delta_{min}=30^\circ$ (see Figure-2). The refractive index of the prism material for blue light is $\sqrt3$. Which of the following statement(s) is(are) correct?
\qfig{p1_phy_s2_q06_1}
\qfig{p1_phy_s2_q06_2}
\begin{qoptions}
\opt{A}{The blue light is polarized in the plane perpendicular to the plane of incidence.}
\opt{B}{The angle of the prism is $45^\circ$.}
\opt{C}{The refractive index of the material of the prism for red light is $\sqrt2$.}
\opt{D}{The angle of refraction for blue light in air at the exit plane of the prism is $60^\circ$.}
\end{qoptions}

\qn{7} In a circuit shown in the figure, the capacitor C is initially uncharged and the key K is open. In this condition, a current of 1 A flows through the 1 $\Omega$ resistor. The key is closed at time $t=t_0$. Which of the following statement(s) is(are) correct? [Given: $e^{-1}=0.36$]
\qfig{p1_phy_s2_q07_1}
\begin{qoptions}
\opt{A}{The value of the resistance R is $\dfrac43\Omega$.}
\opt{B}{For $t<t_0$, the value of current $I_1$ is 3A.}
\opt{C}{At $t=t_0+3.48\ \mu s$, the current in the capacitor is approximately 1 A.}
\opt{D}{For $t\to\infty$, the charge on the capacitor is 12 $\mu C$.}
\end{qoptions}

\sect{Section-I (iii): Matching List Sets}

\qn{8} List-I shows different radioactive decay processes and List-II provides possible emitted particles. Match each entry in List-I with an appropriate entry from List-II, and choose the correct option.
\begin{center}
\begin{tabular}{|c|p{2.5cm}|c|p{3.3cm}|}
\hline
\multicolumn{2}{|c|}{List-I} & \multicolumn{2}{c|}{List-II} \\
\hline
(P) & ${}^{238}_{92}U\to{}^{234}_{91}Pa$ & (1) & one $\alpha$ particle and one $\beta^+$ particle \\
\hline
(Q) & ${}^{214}_{82}Pb\to{}^{210}_{82}Pb$ & (2) & three $\beta^-$ particles and one $\alpha$ particle \\
\hline
(R) & ${}^{210}_{81}Tl\to{}^{206}_{82}Pb$ & (3) & two $\beta^-$ particles and one $\alpha$ particle \\
\hline
(S) & ${}^{228}_{91}Pa\to{}^{224}_{88}Ra$ & (4) & one $\alpha$ particle and one $\beta^-$ particle \\
\hline
 & & (5) & one $\alpha$ particle and two $\beta^+$ particles \\
\hline
\end{tabular}
\end{center}
\begin{qoptions}
\opt{A}{P$\to$4;Q$\to$3;R$\to$2;S$\to$1}
\opt{B}{P$\to$4;Q$\to$1;R$\to$2;S$\to$5}
\opt{C}{P$\to$5;Q$\to$3;R$\to$1;S$\to$4}
\opt{D}{P$\to$5;Q$\to$1;R$\to$3;S$\to$2}
\end{qoptions}

\qn{9} Match the temperature of a black body given in List-I with an appropriate statement in List-II, and choose the correct option. [Given: Wien's constant as $2.9\times10^{-3}$ m-K and $\dfrac{hc}{e}=1.24\times10^{-6}$ V-m]
\begin{center}
\begin{tabular}{|c|c|c|p{4.3cm}|}
\hline
\multicolumn{2}{|c|}{List-I} & \multicolumn{2}{c|}{List-II} \\
\hline
(P) & 2000 K & (1) & The radiation at peak wavelength can lead to emission of photoelectrons from a metal of work function 4 eV \\
\hline
(Q) & 3000 K & (2) & The radiation at peak wavelength is visible to human eye. \\
\hline
(R) & 5000 K & (3) & The radiation at peak emission wavelength will result in the widest central maximum of a single slit diffraction. \\
\hline
(S) & 10000 K & (4) & The power emitted per unit area is 1/16 of that emitted by a blackbody at temperature 6000 K. \\
\hline
 & & (5) & The radiation at peak emission wavelength can be used to image human bones. \\
\hline
\end{tabular}
\end{center}
\begin{qoptions}
\opt{A}{P$\to$3;Q$\to$5;R$\to$2;S$\to$3}
\opt{B}{P$\to$3;Q$\to$2;R$\to$4;S$\to$1}
\opt{C}{P$\to$3;Q$\to$4;R$\to$2;S$\to$1}
\opt{D}{P$\to$1;Q$\to$2;R$\to$5;S$\to$3}
\end{qoptions}

\qn{10} A series LCR circuit is connected to a $100\sin(\omega t)$ Volt source. The resonant angular frequency of the circuit is $10^5$ rad $s^{-1}$ and current amplitude at resonance is $I_0$. When the angular frequency of the source is $\omega=8\times10^4$ rad $s^{-1}$, the current amplitude in the circuit is $0.05I_0$. If $L=100$ mH, match each entry in List-I with an appropriate value from List-II and choose the correct option.
\begin{center}
\begin{tabular}{|c|p{3.5cm}|c|l|}
\hline
\multicolumn{2}{|c|}{List-I} & \multicolumn{2}{c|}{List-II} \\
\hline
(P) & $I_0$ in mA & (1) & 44.4 \\
\hline
(Q) & The quality factor of the circuit & (2) & 4.444 \\
\hline
(R) & The bandwidth of the circuit in rad $s^{-1}$ & (3) & 444.44 \\
\hline
(S) & The peak power dissipated at resonance in Watt & (4) & 2250 \\
\hline
 & & (5) & 0.44 \\
\hline
\end{tabular}
\end{center}
\begin{qoptions}
\opt{A}{P$\to$2;Q$\to$3;R$\to$5;S$\to$1}
\opt{B}{P$\to$3;Q$\to$1;R$\to$4;S$\to$2}
\opt{C}{P$\to$4;Q$\to$5;R$\to$3;S$\to$1}
\opt{D}{P$\to$3;Q$\to$1;R$\to$4;S$\to$1}
\end{qoptions}

\qn{11} A thin conducting rod MN of mass 10 gm, length 25 cm and resistance $5\Omega$ is held on frictionless, long, perfectly conducting vertical rails as shown in the figure. There is a uniform magnetic field $B_0=4$ T directed perpendicular to the plane of the rod-rail arrangement. The rod is released from rest at time t = 0 and it moves down along the rails. Assume air drag is negligible. Match each quantity in List-I with an appropriate value from List-II, and choose the correct option.
[Given: The acceleration due to gravity $g=10\ ms^{-2}$ and $e^{-1}=0.4$]
\qfig{p1_phy_s3_q11_1}
\begin{center}
\begin{tabular}{|c|p{3.5cm}|c|l|}
\hline
\multicolumn{2}{|c|}{List-I} & \multicolumn{2}{c|}{List-II} \\
\hline
(P) & At t = 0.1 s, the magnitude of the induced emf in Volt & (1) & 0.084 \\
\hline
(Q) & At t = 0.1 s, the magnitude of the magnetic force in Newton & (2) & 0.42 \\
\hline
(R) & At t = 0.1 s, the power dissipated as heat in Watt & (3) & 0.042 \\
\hline
(S) & The magnitude of terminal velocity of the rod in m $s^{-1}$ & (4) & 0.035 \\
\hline
 & & (5) & 0.5 \\
\hline
\end{tabular}
\end{center}
\begin{qoptions}
\opt{A}{P$\to$5;Q$\to$2;R$\to$3;S$\to$1}
\opt{B}{P$\to$3;Q$\to$1;R$\to$4;S$\to$5}
\opt{C}{P$\to$2;Q$\to$1;R$\to$4;S$\to$5}
\opt{D}{P$\to$3;Q$\to$4;R$\to$2;S$\to$5}
\end{qoptions}

\sect{Section-II: Numerical Value}

\qn{1} A Hydrogen-like atom has atomic number Z. Photons emitted in the electronic transitions from level n = 3 to level n = 2 in these atoms are used to perform photoelectric effect experiment on a target metal. The maximum kinetic energy of the photoelectrons generated is 3.6 eV. If the photoelectric threshold wavelength for the target metal is 310 nm, the value of Z is \_\_\_\_\_\_\_\_. [Given: hc = 1240 eV-nm and Rhc = 13.6 eV, where R is the Rydberg constant, h is the Planck's constant and c is the speed of light in vacuum]

\qn{2} An optical arrangement consists of two concave mirrors $M_1$ and $M_2$, and a convex lens L with a common principal axis, as shown in the figure. The focal length of L is 10 cm. The radii of curvature of $M_1$ and $M_2$ are 20 cm and 24 cm, respectively. The distance between L and $M_2$ is 20 cm. A point object S is placed at the mid-point between L and $M_2$ on the axis. When the distance between L and $M_1$ is $n/7$ cm, the image formed after refraction from lens then reflection from $M_1$ and again refraction from lens and finally reflection from $M_2$ coincides with S. The value of n is \_\_\_\_\_\_\_\_.
\qfig{p1_phy_s4_q02_1}

\qn{3} In an experiment for determination of the focal length of a thin convex lens, the distance of the object from the lens is $10\pm0.1$ cm and the distance of its real image from the lens is $10\pm0.2$ cm. The error in the determination of focal length of the lens is $n\%$. The value of 100n is \_\_\_\_\_\_\_\_.

\qn{4} A closed container contains a homogeneous mixture of two moles of an ideal gas ($\gamma=4/3$) and one mole of an ideal gas ($\gamma=9/8$). Here, $\gamma$ is the ratio of the specific heats at constant pressure and constant volume of an ideal gas. The gas mixture does a work of 66 Joule when heated at constant pressure. The change in its internal energy is \_\_\_\_\_\_\_\_ Joule.

\qn{5} A person of height 1.6 m is walking away from a lamp post of height 4 m along a straight path on the flat ground. The lamp post and the person are always perpendicular to the ground. If the speed of the person is 60 cm $s^{-1}$, the speed of the tip of the person's shadow on the ground with respect to the ground is \_\_\_\_\_\_\_\_ cm $s^{-1}$.

\qn{6} Two point-like objects of masses 20 gm and 20 gm are fixed at the two ends of a rigid massless rod of length 20 cm. This system is suspended vertically from a rigid ceiling using a thin wire attached to its center of mass, as shown in the figure. The resulting torsional pendulum undergoes small oscillations. The torsional constant of the wire is $1.6\times10^{-7}$ N m rad$^{-1}$. The angular frequency of the oscillations is $n\times10^{-2}$ rad $s^{-1}$. The value of n is \_\_\_\_\_.
\qfig{p1_phy_s4_q06_1}

\partsec{2}{CHEMISTRY}
\sect{Section-I (i): Single Correct}

\qn{1} In the scheme given below which of the following is correct for X and Y.
\qfig{p1_chem_s1_q01_1}
\begin{qoptions}
\opt{A}{X is a powerful reducing agent}
\opt{B}{Y is a colorless odourless gas}
\opt{C}{Color in X is due to Ligand-Metal charge transfer}
\opt{D}{Y contains both $\sigma$ and $\pi$ bonds}
\end{qoptions}

\qn{2} Plotting $1/\Lambda_m$ against $c\Lambda_m$ for aqueous solutions of a monobasic weak acid (HX) resulted in a straight line with y-axis intercept of P and slope of S. The ratio P/S is\\
\relax[$\Lambda_m$ = Molar conductivity of HX at concentration 'C'.\\
$\Lambda_m^o$ = Limiting molar conductivity of HX.\\
c = Molar concentration of HX.\\
$K_a$ = Dissociation constant of HX]
\begin{qoptions}
\opt{A}{$K_a\Lambda_m^o$}
\opt{B}{$K_a\Lambda_m^o/2$}
\opt{C}{$2K_a\Lambda_m^o$}
\opt{D}{$1/(K_a\Lambda_m^o)$}
\end{qoptions}

\qn{3} On decreasing the pH from 7 to 4, the solubility of a sparingly soluble salt (MX) of a weak acid (HX) increased from $10^{-4}$ mol $L^{-1}$ to $10^{-3}$ mol $L^{-1}$. The $pK_a$ of HX is:
\begin{qoptions}
\opt{A}{3}
\opt{B}{4}
\opt{C}{6}
\opt{D}{2}
\end{qoptions}

\qn{4} In the given reaction scheme, P is a phenyl alkyl ether, Q is an aromatic compound; R and S are the major products.
$$P\xrightarrow{HI}Q\xrightarrow[\text{(iii)}H_3O^+]{\text{(i)}NaOH\text{ (ii)}CO_2}R\xrightarrow[\text{(ii)}H_3O^+]{\text{(i) }CH_3-C(=O)-Cl}S$$
The correct statement about S is
\begin{qoptions}
\opt{A}{It primarily inhibits noradrenaline degrading enzymes.}
\opt{B}{It increase the synthesis of prostaglandin.}
\opt{C}{It is potent vasodialator, analgesic and anti inflamantory}
\opt{D}{It is \textit{ortho}-acetylbenzoic acid.}
\end{qoptions}

\sect{Section-I (ii): One or More Correct}

\qn{5} The correct statement(s) related to processes involved in the extraction of metals is(are)
\begin{qoptions}
\opt{A}{Roasting of Chalcocite produces Malachite.}
\opt{B}{Calcination of Dolomite produces Quick Lime as one of the product.}
\opt{C}{Copper pyrites is heated with silica in a Electric arc furnace to remove iron.}
\opt{D}{Impure gold is treated with aqueous NaCN in the presence of oxygen followed by reduction with zinc metal.}
\end{qoptions}

\qn{6} In the following reactions, P, Q, R, and S are the major products.\\
$CH_3CH_2CH(CH_3)CH_2CN\xrightarrow[\text{(ii)}CH_3MgBr,\text{ then }H_2O]{\text{(i)}CH_3MgBr,\text{ then }H_3O^\oplus}P$
\qfig{p1_chem_s2_q06_1}
\qfig{p1_chem_s2_q06_2}
\qfig{p1_chem_s2_q06_3}
The correct statement(s) about P, Q, R, and S is(are)
\begin{qoptions}
\opt{A}{Both P and R have asymmetric carbon(s).}
\opt{B}{P has asymmetric carbon(s), S does \textbf{not} have any asymmetric carbon.}
\opt{C}{Both P and Q have asymmetric carbon(s).}
\opt{D}{Both Q and R have asymmetric carbon(s).}
\end{qoptions}

\qn{7} Consider the following reaction scheme and choose the correct option(s) for the major products Q, R and S.\\
\resizebox{\linewidth}{!}{$\text{Styrene}\xrightarrow[\text{(ii)}NaOH,H_2O_2,H_2O]{\text{(i)}B_2H_6}P\xrightarrow[\text{(iii)}H_2O]{\text{(i)}CrO_3,H_2SO_4\text{ (ii)}Cl_2,\text{ Red Phosphorus}}Q$}\\
$$P\xrightarrow[\text{(ii)}NaCN\text{ (iii)}H_3O^+,\Delta]{\text{(i)}SOCl_2}R\xrightarrow[\text{(ii)}Zn-Hg/HCl]{\text{(i) conc.}H_2SO_4}S$$
\begin{qoptions}
\opt{A}{\qfigi{p1_chem_s2_q07_A}}
\opt{B}{\qfigi{p1_chem_s2_q07_B}}
\opt{C}{\qfigi{p1_chem_s2_q07_C}}
\opt{D}{\qfigi{p1_chem_s2_q07_D}}
\end{qoptions}

\sect{Section-I (iii): Matching List Sets}

\qn{8} Match the reactions (in the given stoichiometry of the reactants) in List-I with one of their products given in List-II and choose the correct option.
\begin{center}
\begin{tabular}{|c|p{3.15cm}|c|p{2.7cm}|}
\hline
\multicolumn{2}{|c|}{List-I} & \multicolumn{2}{c|}{List-II} \\
\hline
(P) & $P_2O_3+3H_2O\to$ & (1) & $P(O)(OCH_3)Cl_2$ \\
\hline
(Q) & $P_4+3NaOH+3H_2O\to$ & (2) & $H_3PO_3$ \\
\hline
(R) & $PCl_5+CH_3COOH\to$ & (3) & $PH_3$ \\
\hline
(S) & $H_3PO_2+2H_2O+4AgNO_3\to$ & (4) & $POCl_3$ \\
\hline
 & & (5) & $H_3PO_4$ \\
\hline
\end{tabular}
\end{center}
\begin{qoptions}
\opt{A}{P$\to$2;Q$\to$3;R$\to$4;S$\to$5}
\opt{B}{P$\to$3;Q$\to$5;R$\to$4;S$\to$2}
\opt{C}{P$\to$5;Q$\to$2;R$\to$1;S$\to$3}
\opt{D}{P$\to$2;Q$\to$3;R$\to$1;S$\to$5}
\end{qoptions}

\qn{9} Match the electronic configurations in List-I with appropriate metal complex ions in List-II and choose the correct option.\\
\relax[Atomic Number: Fe = 26, Mn = 25, Co = 27]
\begin{center}
\begin{tabular}{|c|c|c|l|}
\hline
\multicolumn{2}{|c|}{List-I} & \multicolumn{2}{c|}{List-II} \\
\hline
(P) & $t_{2g}^6e_g^0$ & (1) & $[Fe(H_2O)_6]^{2+}$ \\
\hline
(Q) & $t_{2g}^3e_g^2$ & (2) & $[Mn(H_2O)_6]^{2+}$ \\
\hline
(R) & $e^4t_2^3$ & (3) & $[Co(NH_3)_6]^{3+}$ \\
\hline
(S) & $t_{2g}^4e_g^2$ & (4) & $[FeCl_4]^-$ \\
\hline
 & & (5) & $[CoCl_4]^{2-}$ \\
\hline
\end{tabular}
\end{center}
\begin{qoptions}
\opt{A}{P$\to$1;Q$\to$4;R$\to$2;S$\to$3}
\opt{B}{P$\to$1;Q$\to$2;R$\to$4;S$\to$5}
\opt{C}{P$\to$3;Q$\to$2;R$\to$5;S$\to$1}
\opt{D}{P$\to$3;Q$\to$2;R$\to$4;S$\to$1}
\end{qoptions}

\qn{10} Match the reactions in List-I with the features of their products in List-II and choose the correct option.
\begin{center}
\begin{tabular}{|c|p{3.6cm}|c|p{2.2cm}|}
\hline
\multicolumn{2}{|c|}{List-I} & \multicolumn{2}{c|}{List-II} \\
\hline
(P) & (-)-1-chloro-2-ethylpentane (single enantiomer)\newline \resizebox{2.6cm}{!}{$\xrightarrow[S_N2\text{ reaction}]{aq.\ NaOH}$} & (1) & Retention of configuration \\
\hline
(Q) & (-)-2-chloropentane (single enantiomer)\newline \resizebox{2.6cm}{!}{$\xrightarrow[S_N2\text{ reaction}]{aq.\ NaOH}$} & (2) & Mixture of diastereomers \\
\hline
(R) & (-)-3-chloro-3-methylhexane (single enantiomer)\newline \resizebox{2.6cm}{!}{$\xrightarrow[S_N1\text{ reaction}]{aq.\ NaOH}$} & (3) & Inversion of configuration \\
\hline
(S) & \qfigi{p1_chem_s3_q10_4} (Single enantiomer)\newline \resizebox{2.6cm}{!}{$\xrightarrow[S_N1\text{ reaction}]{aq.\ NaOH}$} & (4) & Mixture of enantiomers \\
\hline
 & & (5) & Mixture of structural isomers \\
\hline
\end{tabular}
\end{center}
\begin{qoptions}
\opt{A}{P$\to$1;Q$\to$2;R$\to$5;S$\to$3}
\opt{B}{P$\to$2;Q$\to$1;R$\to$3;S$\to$5}
\opt{C}{P$\to$1;Q$\to$2;R$\to$5;S$\to$4}
\opt{D}{P$\to$1;Q$\to$3;R$\to$4;S$\to$2}
\end{qoptions}

\qn{11} The major products obtained from the reactions in List-II are the reactants for the named reactions mentioned in List-I. Match List-I with List-II and choose the correct option.
\begin{center}
\begin{tabular}{|c|p{2.3cm}|c|p{3.4cm}|}
\hline
\multicolumn{2}{|c|}{List-I} & \multicolumn{2}{c|}{List-II} \\
\hline
(P) & Cumene hydroperoxide rearrangement & (1) & Acetophenone $\xrightarrow{Zn-Hg,\ HCl}$ \\
\hline
(Q) & Gattermann reaction & (2) & Toluene\newline \resizebox{3.2cm}{!}{$\xrightarrow[\text{(ii)}SOCl_2]{\text{(i)}KMnO_4,KOH,\Delta}$} \\
\hline
(R) & Gattermann-Koch reaction & (3) & Benzene\newline \resizebox{3.2cm}{!}{$\xrightarrow[anhyd.\ AlCl_3]{CH_3CH_2CH_2Cl}$} \\
\hline
(S) & Rosenmund reduction & (4) & Aniline $\xrightarrow[273-278\ K]{NaNO_2/HCl}$ \\
\hline
 & & (5) & Phenol $\xrightarrow{Zn,\ \Delta}$ \\
\hline
\end{tabular}
\end{center}
\begin{qoptions}
\opt{A}{P$\to$3;Q$\to$4;R$\to$5;S$\to$2}
\opt{B}{P$\to$1;Q$\to$3;R$\to$5;S$\to$2}
\opt{C}{P$\to$2;Q$\to$3;R$\to$1;S$\to$4}
\opt{D}{P$\to$2;Q$\to$4;R$\to$1;S$\to$3}
\end{qoptions}

\sect{Section-II: Numerical Value}

\qn{1} The stoichiometric reaction of 516 g of dimethyldichlorosilane with water results in a tetrameric cyclic product X in 25\% yield. The weight (in g) of X obtained is\_\_\_. [Use, molar mass (g $mol^{-1}$): H = 1, C = 12, O = 16, Si = 28, Cl = 35.5]

\qn{2} A gas has a compressibility factor of 0.6 and a molar volume of 0.4 $dm^3\,mol^{-1}$ at a temperature of 800 K and pressure x atm. If it shows ideal gas behaviour at the same temperature and pressure, the molar volume will be y $dm^3\,mol^{-1}$. The value of x/y is \_\_\_. [Use: Gas constant, $R=8\times10^{-2}$ L atm $K^{-1}mol^{-1}$]

\qn{3} The plot of $\log k_f$ versus $1/T$ for a reversible reaction $A(g)\rightleftharpoons P(g)$ is shown.
\qfig{p1_chem_s4_q03_1}
Pre-exponential factors for the forward and backward reactions are $10^{15}\ s^{-1}$ and $10^{11}\ s^{-1}$, respectively. If the value of log K for the reaction at 500 K is 5, the value of $|\log k_b|$ at 250 K is \_\_\_\_\_\_. [K = equilibrium constant of the reaction, $k_f$ = rate constant of forward reaction, $k_b$ = rate constant of backward reaction]

\qn{4} One mole of an ideal monoatomic gas undergoes two reversible processes ($A\to B$ and $B\to C$) as shown in the given figure:
\qfig{p1_chem_s4_q04_1}
$A\to B$ is an adiabatic process. If the total heat absorbed in the entire process ($A\to B$ and $B\to C$) is $RT_2\ln10$, the value of $2\log V_3$ is \_\_\_\_\_\_\_\_. [Use, molar heat capacity of the gas at constant pressure, $C_{p,m}=\dfrac52R$]

\qn{5} In an one-litre flask, 6 moles of A undergoes the reaction $A(g)\rightleftharpoons P(g)$. The progress of product formation at two temperatures (in Kelvin), $T_1$ and $T_2$, is shown in the figure:
\qfig{p1_chem_s4_q05_1}
If $T_1=2T_2$ and $(\Delta G_2^\ominus-\Delta G_1^\ominus)=RT_2\ln x$, then the value of x is \_\_\_\_\_\_. [$\Delta G_1^\ominus$ and $\Delta G_2^\ominus$ are standard Gibb's free energy change for the reaction at temperatures $T_1$ and $T_2$, respectively.]

\qn{6} The total number of $sp^2$ hybridised carbon atoms in the major product P (a non-heterocyclic compound) of the following reaction is \_\_\_\_\_\_\_.
\qfig{p1_chem_s4_q06_1}

\partsec{3}{MATHEMATICS}
\sect{Section-I (i): Single Correct}

\qn{1} Let $f:(0,1)\to R$ and $g:(0,1)\to R$ be the functions defined as $f(x)=\sqrt n$; $x\in\left[\dfrac{1}{n+1},\dfrac1n\right)$; $n\in N$,\\
\resizebox{\linewidth}{!}{$\sqrt2\displaystyle\int_{x^2}^x\sqrt{\dfrac{1-t}{t}}\,dt<g(x)<2\sqrt{2x-3x^2+4x^3-3x^4}\ \forall x\in(0,1).$}\\
Then $\displaystyle\lim_{x\to0}(f(x)\cdot g(x))$
\begin{qoptions}
\opt{A}{does \textbf{NOT} exist}
\opt{B}{equals $\sqrt2$}
\opt{C}{equals $2\sqrt2$}
\opt{D}{equals $3\sqrt2$}
\end{qoptions}

\qn{2} Let C be the cube with set of vertices $\{(x,y,z)\in R^3:x,y,z\in\{-1,1\}\}$. Let $D_f$ be the set of all twelve lines containing the diagonals of the six faces of the cube C. Let $D_m$ be the set of all four lines containing the main diagonals of the cube C; for instance, the line passing through the vertices $(-1,-1,-1)$ and $(1,1,1)$ is in $D_m$. Let $d(\ell_1,\ell_2)$ denotes the shortest distance between lines $\ell_1$ and $\ell_2$. If the maximum value of $d(\ell_1,\ell_2)$, as $\ell_1$ varies over $D_f$ and $\ell_2$ varies over $D_m$ is equal to $\lambda$, then $\lambda^{-2}$ equals
\begin{qoptions}
\opt{A}{2}
\opt{B}{6}
\opt{C}{3}
\opt{D}{1.5}
\end{qoptions}

\qn{3} Let P,Q,R be any three points chosen randomly from
$$X=\left\{(x,y)\in Z^2:\left(\dfrac{x-2}{3}\right)^2+\dfrac{y^2}{20}<1<\dfrac{5x-10}{y^2}\right\}.$$
If p denotes the probability that area of $\triangle PQR$ is a natural number, then 220p equals
\begin{qoptions}
\opt{A}{71}
\opt{B}{72}
\opt{C}{73}
\opt{D}{74}
\end{qoptions}

\qn{4} Let A be a point lying on the parabola $P\equiv y^2=4ax$; $a\in N$. The normal to parabola P at point A meets X-axis at point B. The area of $\triangle FAB$, where F is focus of parabola, is 120 square units. If slope m of normal at point A is a natural number, then $2a+3m$ equals
\begin{qoptions}
\opt{A}{13}
\opt{B}{12}
\opt{C}{14}
\opt{D}{11}
\end{qoptions}

\sect{Section-I (ii): One or More Correct}

\qn{5} Let $S=\left\{x:\dfrac{x^3-9x^2+23x-15}{x^3-12x^2+44x-48}<0\right\}$; $T=\left\{\dfrac{x}{y}:x,y\in\{1,2,3\}\right\}$.\\
Then which of the following statement is(are) true?
\begin{qoptions}
\opt{A}{There are 35 increasing functions from S to T.}
\opt{B}{There are more than 350 continuous functions from S to T.}
\opt{C}{There are less than 350 differentiable functions from S to T.}
\opt{D}{Every continuous function from S to T is differentiable}
\end{qoptions}

\qn{6} Let $T_1$ and $T_2$ be two distinct common tangents to the ellipse $E\equiv x^2+2y^2+3=2x+8y$ and parabola $P\equiv y^2+16=4y+12x$. Suppose $T_1$ touches E and P at $E_1$ and $P_1$ respectively. Suppose $T_2$ touches E and P at $E_2$ and $P_2$ respectively. If slope of $T_1$ is greater than that of $T_2$, then which of the following options is/are correct?
\begin{qoptions}
\opt{A}{The area of quadrilateral $E_1P_1P_2E_2$ equals 35 square units.}
\opt{B}{Point of intersection of $T_1$ and $T_2$ is at 5 units from $(1,-2)$.}
\opt{C}{$E_1P_2+E_2P_1=2\sqrt{74}$}
\opt{D}{The area of $\triangle E_1P_1E_2$ equals 5 square units.}
\end{qoptions}

\qn{7} Consider the square region $S=[0,1]\times[0,1]$.\\
Let $G=\{(x,y)\in S:y>f(x)\}$ be called the GREEN region\\
$R=\{(x,y)\in S:y<f(x)\}$ be called the RED region.\\
$L_h=\{(x,h)\in S:x\in[0,1]\}$ be the horizontal line drawn at a height $h\in[0,1]$.\\
If $f:[0,1]\to[0,1]$ be the function defined as
$$f(x)=\dfrac{x^3}{3}-x^2+\dfrac{5x}{9}+\dfrac{11}{36},$$
then which of the following statements is(are) true?
\begin{qoptions}
\opt{A}{There exists an $h\in\left[\dfrac16,\dfrac23\right]$ such that area of GREEN region above line $L_h$ equals that of GREEN region below line $L_h$.}
\opt{B}{There exists an $h\in\left[\dfrac16,\dfrac23\right]$ such that area of RED region above line $L_h$ equals that of RED region below line $L_h$.}
\opt{C}{There exists an $h\in\left[\dfrac16,\dfrac23\right]$ such that area of GREEN region above line $L_h$ equals that of RED region below line $L_h$.}
\opt{D}{There exists an $h\in\left[\dfrac16,\dfrac23\right]$ such that area of RED region above line $L_h$ equals that of GREEN region below line $L_h$.}
\end{qoptions}

\sect{Section-I (iii): Matching List Sets}

\qn{8} Let $\alpha,\beta$ and $\gamma$ be real numbers. Consider the following system of linear equations\\
$x+2y+z=7$\\
$x+2\alpha z=11$\\
$2x-3y+3\beta z=4\gamma$\\
Match each entry in List-I to the correct entries in List-II
\begin{center}
\begin{tabular}{|c|p{3.6cm}|c|p{2.3cm}|}
\hline
\multicolumn{2}{|c|}{List-I} & \multicolumn{2}{c|}{List-II} \\
\hline
(P) & If $3\beta=\dfrac12(14\alpha-3)$ and $\gamma=7$, then the system has & (1) & a unique solution \\
\hline
(Q) & If $3\beta=\dfrac12(14\alpha-3)$ and $\gamma\ne7$, then the system has & (2) & no solution \\
\hline
(R) & If $3\beta\ne\dfrac12(14\alpha-3)$ where $2\alpha=1$ and $\gamma\ne7$, then the system has & (3) & infinitely many solutions \\
\hline
(S) & If $3\beta\ne\dfrac12(14\alpha-3)$ where $2\alpha=1$ and $\gamma=7$, then the system has & (4) & x = 11, y = --2 and z = 0 as a solution \\
\hline
 & & (5) & x = --15, y = 4 and z = 0 as a solution \\
\hline
\end{tabular}
\end{center}
The correct option is:
\begin{qoptions}
\opt{A}{P$\to$2;Q$\to$1;R$\to$1;S$\to$3}
\opt{B}{P$\to$3;Q$\to$2;R$\to$1;S$\to$4}
\opt{C}{P$\to$3;Q$\to$2;R$\to$5;S$\to$4}
\opt{D}{P$\to$2;Q$\to$1;R$\to$4;S$\to$5}
\end{qoptions}

\qn{9} Consider the given data with frequency distribution\\
$x_i$: 3 8 11 10 5 4\\
$f_i$: 5 2 3 2 4 4\\
Match each entry in List-I to the correct entries in List-II.
\begin{center}
\begin{tabular}{|c|p{4.3cm}|c|l|}
\hline
\multicolumn{2}{|c|}{List-I} & \multicolumn{2}{c|}{List-II} \\
\hline
(P) & The mean of the above data is & (1) & 2.5 \\
\hline
(Q) & The median of the above data is & (2) & 2.4 \\
\hline
(R) & The mean deviation about the mean of the above data is & (3) & 2.7 \\
\hline
(S) & The mean deviation about the median of the above data is & (4) & 6 \\
\hline
 & & (5) & 5 \\
\hline
\end{tabular}
\end{center}
The correct option is:
\begin{qoptions}
\opt{A}{P$\to$4;Q$\to$5;R$\to$3;S$\to$2}
\opt{B}{P$\to$4;Q$\to$5;R$\to$1;S$\to$2}
\opt{C}{P$\to$5;Q$\to$4;R$\to$3;S$\to$1}
\opt{D}{P$\to$4;Q$\to$4;R$\to$2;S$\to$2}
\end{qoptions}

\qn{10} Let $\ell_1$ and $\ell_2$ be the lines $\vec r_1=\lambda(\hat i+\hat j+\hat k)$ and $\vec r_2=(\hat j-\hat k)+\mu(\hat i+\hat k)$, respectively. Let X be the set of all the planes H that contain the line $\ell_1$. For a plane H, let d(H) denote the smallest possible distance between the points of $\ell_2$ and H. Let $H_0$ be plane in X for which $d(H_0)$ is the maximum value of d(H) as H varies over all planes in X.\\
Match each entry in List-I to the correct entries in List-II.
\begin{center}
\begin{tabular}{|c|p{4.3cm}|c|l|}
\hline
\multicolumn{2}{|c|}{List-I} & \multicolumn{2}{c|}{List-II} \\
\hline
(P) & The value of $d(H_0)$ is & (1) & 0 \\
\hline
(Q) & The distance of the point (0,1,2) from $H_0$ is & (2) & $\dfrac{1}{\sqrt3}$ \\
\hline
(R) & The distance of origin from $H_0$ is & (3) & $\sqrt3$ \\
\hline
(S) & The distance of origin from the point of intersection of planes y = z, x = 1 and $H_0$ is & (4) & $\sqrt2$ \\
\hline
 & & (5) & $\dfrac{1}{\sqrt2}$ \\
\hline
\end{tabular}
\end{center}
The correct option is:
\begin{qoptions}
\opt{A}{P$\to$2;Q$\to$4;R$\to$5;S$\to$3}
\opt{B}{P$\to$5;Q$\to$4;R$\to$1;S$\to$3}
\opt{C}{P$\to$2;Q$\to$3;R$\to$1;S$\to$2}
\opt{D}{P$\to$5;Q$\to$3;R$\to$4;S$\to$2}
\end{qoptions}

\qn{11} Let z be complex number satisfying $|z|^3+2z^2+4\bar z-8=0$, where $\bar z$ denotes the complex conjugate of z. Let the imaginary part of z be nonzero.\\
Match each entry in List-I to the correct entries in List-II.
\begin{center}
\begin{tabular}{|c|c|c|l|}
\hline
\multicolumn{2}{|c|}{List-I} & \multicolumn{2}{c|}{List-II} \\
\hline
(P) & $|z|^2$ is equal to & (1) & 12 \\
\hline
(Q) & $|z-\bar z|^2$ is equal to & (2) & 4 \\
\hline
(R) & $|z|^2+|z+\bar z|^2$ is equal to & (3) & 8 \\
\hline
(S) & $|z+1|^2$ is equal to & (4) & 10 \\
\hline
 & & (5) & 7 \\
\hline
\end{tabular}
\end{center}
The correct option is:
\begin{qoptions}
\opt{A}{P$\to$2;Q$\to$3;R$\to$5;S$\to$4}
\opt{B}{P$\to$1;Q$\to$3;R$\to$5;S$\to$4}
\opt{C}{P$\to$2;Q$\to$1;R$\to$3;S$\to$5}
\opt{D}{P$\to$2;Q$\to$4;R$\to$5;S$\to$1}
\end{qoptions}

\sect{Section-II: Numerical Value}

\qn{1} Let $2\tan^{-1}x\in(-\pi,\pi)\ \forall x\in R$. The number of real solutions of\\
\resizebox{\linewidth}{!}{$\sqrt{1+\cos\dfrac{2x}{3}}=\sqrt2\tan^{-1}\tan\dfrac x3;\ x\in\left(-\dfrac{9\pi}{2},\dfrac{9\pi}{2}\right);\ x\ne\dfrac{(6n+3)\pi}{2},\ n\in Z$}\\
is equal to

\qn{2} Let $n\in N$, $n>1$ and $f:[0,1]\to R$ be the function defined by
$$f(x)=\begin{cases}n(1-2nx), & 0\le x\le\dfrac{1}{2n},\\[4pt] 2n(2nx-1), & \dfrac{1}{2n}\le x\le\dfrac{3}{4n},\\[4pt] 4n(1-nx), & \dfrac{3}{4n}\le x\le\dfrac1n,\\[4pt] \dfrac{n}{n-1}(nx-1), & \dfrac1n\le x\le1.\end{cases}$$
If n is such that the area of the region bounded by the curves x = 0, x = 1, y = 0 and y = f(x) is 3 square units, then the maximum value of the function f is

\qn{3} Let $a_r$ be an (r+2) digit natural number where first and last digits are 5, and the remaining digits are 7 i.e. $a_r=5\underbrace{77\ldots7}_{r}5$. If $\displaystyle\sum_{r=0}^{98}a_r=\dfrac{a_{99}-m}{n}$; m,n$\in N$; m,n < 3000, then value of m + n equals.

\qn{4} If set A defined as\\
\resizebox{\linewidth}{!}{$A=\left\{\dfrac{a+bi}{c+di}\in C:\begin{pmatrix}a&b&1\\c&d&1\\1&1&1\end{pmatrix}\begin{pmatrix}0&1&0\\1&0&0\\0&0&1\end{pmatrix}=\begin{pmatrix}1686\sin\theta&1967&1\\-3\cos\theta&7&1\\1&1&1\end{pmatrix},\ \theta\in R\right\}$}\\
contains only one positive integer n, then n equals

\qn{5} Let\\
\resizebox{\linewidth}{!}{$A=\left\{x\hat i+y\hat j+z\hat k:(x,y,z)\in R^3,x^2+y^2+z^2=1\right\}$,}\\
\resizebox{\linewidth}{!}{$B=\left\{x\hat i+y\hat j+z\hat k:(x,y,z)\in R^3,\left|\sqrt3x+2y+3z-16\right|=14\right\}$.}\\
Let $\vec a,\vec b,\vec c$ be three distinct vectors in set $A\cap B$ such that $|\vec a-\vec b|=|\vec b-\vec c|=|\vec c-\vec a|$. If V denotes the volume of parallelopiped determined by the vectors $\vec a,\vec b,\vec c$ then $16\sqrt3V$ equals

\qn{6} Let a,b$\in R-\{0\}$. If coefficient of $x^5$ in binomial expansion of $\left(\dfrac{ax^2}{4}+\dfrac{140}{27bx}\right)^4$ is equal to coefficient of $x^{-5}$ in $\left(2ax-\dfrac{1}{4bx^2}\right)^7$, then 2b is equal to

\papersec{2}
\partsec{1}{PHYSICS}
\sect{Section-I (i): Single Correct}

\qn{1} An electric dipole is formed by two charges $+q$ and $-q$ located in xy-plane at (0, 2) mm and $(0,-2)$ mm, respectively, as shown in the figure. The electric potential at point P(100, 100) mm due to the dipole is $V_0$. The charges $+q$ and $-q$ are then moved to the points $(-2,2)$ mm and $(2,-2)$ mm, respectively. What is the value of electric potential at P due to the new dipole?
\qfig{p2_phy_s1_q01_1}
\begin{qoptions}
\opt{A}{$V_0/4$}
\opt{B}{$V_0/2$}
\opt{C}{$V_0/\sqrt2$}
\opt{D}{Zero}
\end{qoptions}

\qn{2} Young's modulus of elasticity Y is expressed in terms of three derived quantities, namely, the gravitational constant G, Planck's constant h and the speed of light c, as $Y=c^\alpha h^\beta G^\gamma$. Which of the following is the correct option?
\begin{qoptions}
\opt{A}{$\alpha=7,\beta=-1,\gamma=-2$}
\opt{B}{$\alpha=-7,\beta=-1,\gamma=-2$}
\opt{C}{$\alpha=7,\beta=-1,\gamma=2$}
\opt{D}{$\alpha=-7,\beta=1,\gamma=-2$}
\end{qoptions}

\qn{3} A particle of mass m is moving in the xy-plane such that its velocity at a point (x, y) is given as $\vec v=\alpha\left(y\hat x+2x^2\hat y\right)$, where $\alpha$ is a non-zero constant. What is the force $\vec F$ acting on the particle?
\begin{qoptions}
\opt{A}{$\vec F=2m\alpha^2(x\hat x+y\hat y)$}
\opt{B}{$\vec F=m\alpha^2(y\hat x+2x\hat y)$}
\opt{C}{$\vec F=2m\alpha^2(y\hat x+x\hat y)$}
\opt{D}{$\vec F=2m\alpha^2(x^2\hat x+2xy\hat y)$}
\end{qoptions}

\qn{4} An ideal gas is in thermodynamic equilibrium. The number of degrees of freedom of a molecule of the gas in n. The internal energy of one mole of the gas is $U_n$ and the speed of sound in the gas is $v_n$. At a fixed temperature and pressure, which of the following is the correct option?
\begin{qoptions}
\opt{A}{$v_3<v_6$ and $U_3>U_6$}
\opt{B}{$v_5>v_3$ and $U_3>U_5$}
\opt{C}{$v_5>v_7$ and $U_5<U_7$}
\opt{D}{$v_6<v_7$ and $U_6<U_7$}
\end{qoptions}

\sect{Section-I (ii): One or More Correct}

\qn{5} A monochromatic light wave is incident normally on a glass slab of thickness d, as shown in the figure. The refractive index of the slab increases linearly from $n_1$ to $n_2$ over the height h. Which of the following statement(s) is (are) true about the light wave emerging out of the slab?
\qfig{p2_phy_s2_q05_1}
\begin{qoptions}
\opt{A}{It will deflect up by an angle $\tan^{-1}\left[\dfrac{(n_2^2-n_1^2)d}{2h}\right]$}
\opt{B}{It will deflect up by an angle $\tan^{-1}\left[\dfrac{(n_2-n_1)d}{h}\right]$}
\opt{C}{It will not deflect.}
\opt{D}{The deflection angle depends only on $(n_2-n_1)$ and not on the individual values of $n_1$ and $n_2$.}
\end{qoptions}

\qn{6} An annular disk of mass M, inner radius a and outer radius b is placed on a horizontal surface with coefficient of friction $\mu$, as shown in the figure. At some time, an impulse $J_0\hat x$ is applied at a height h above the center of the disk. If $h=h_m$ then the disk rolls without slipping along the x-axis. Which of the following statement(s) is(are) correct?
\qfig{p2_phy_s2_q06_1}
\begin{qoptions}
\opt{A}{For $\mu\ne0$ and $a\to0$, $h_m=b/2$}
\opt{B}{For $\mu\ne0$ and $a\to b$, $h_m=b$}
\opt{C}{For $h=h_m$, the initial angular velocity depends on the inner radius a.}
\opt{D}{For $\mu=0$ and $h=0$, the wheel always slides without rolling.}
\end{qoptions}

\qn{7} The electric field associated with an electromagnetic wave propagating in a dielectric medium is given by $\vec E=30(2\hat x+\hat y)\sin\left[2\pi\left(5\times10^{14}t-\dfrac{10^7}{3}z\right)\right]$ V $m^{-1}$. Which of the following option(s) is(are) correct? [Given: The speed of light in vacuum, $c=3\times10^8\,ms^{-1}$]
\begin{qoptions}
\opt{A}{$B_x=-2\times10^{-7}\sin\left[2\pi\left(5\times10^{14}t-\dfrac{10^7}{3}z\right)\right]$ Wb$m^{-2}$}
\opt{B}{$B_y=2\times10^{-7}\sin\left[2\pi\left(5\times10^{14}t-\dfrac{10^7}{3}z\right)\right]$ Wb$m^{-2}$}
\opt{C}{The wave is polarized in the xy-plane with polarization angle $\tan^{-1}\left(\dfrac12\right)$ with respect to the x-axis.}
\opt{D}{The refractive index of the medium is 1.5.}
\end{qoptions}

\sect{Section-II (i): Paragraph Type}

\paraq{1}{2}
$S_1$ and $S_2$ are two identical sound sources of frequency 656 Hz. The source $S_1$ is located at O and $S_2$ moves anti-clockwise with a uniform speed $8\sqrt2\ ms^{-1}$ on a circular path around O, as shown in the figure. There are three points P, Q and R on this path such that P and R are diametrically opposite while Q is equidistant from them. A sound detector is placed at point P. The source $S_1$ can move along direction OP. [Given: The speed of sound in air is 324 $ms^{-1}$]
\qfig{p2_phy_s3_para1_2_1}

\qn{1} When only $S_2$ is emitting sound and it is at Q, the frequency of sound measured by the detector in Hz is \_\_\_\_\_\_\_\_\_.

\qn{2} Consider both sources emitting sound. When $S_2$ is at R and $S_1$ approaches the detector with a speed 8 $ms^{-1}$, the beat frequency measured by the detector is \_\_\_\_\_\_\_Hz.

\paraq{3}{4}
A cylindrical furnace has height (H) and diameter (D) both 1 m. It is maintained at temperature 360 K. The air gets heated inside the furnace at constant pressure $P_a$ and its temperature becomes T = 360 K. The hot air with density $\rho$ rises up a vertical chimney of diameter d = 0.1 m and height h = 8 m above the furnace and exits the chimney (see the figure). As a result, atmospheric air of density $\rho_a=1.2$ kg $m^{-3}$, pressure $P_a$ and temperature $T_a=$ 300 K enters the furnace. Assume air as an ideal gas, neglect the variations in $\rho$ and T inside the chimney and the furnace. Also ignore the viscous effects. Assume atmospheric pressure $P_a$ at base of furnace outside the structure. [Given: The acceleration due to gravity $g=10\ ms^{-2}$ and $\pi=3.14$]
\qfig{p2_phy_s3_para3_4_1}

\qn{3} Considering the air flow to be streamline, the steady velocity of air exiting the chimney is \_\_\_\_\_\_\_\_ m $s^{-1}$.

\qn{4} When the chimney is closed using a cap at the top, a pressure difference $\Delta P$ develops between the top and the bottom surfaces of the cap. If the changes in the temperature and density of the hot air, due to the stoppage of air flow, are negligible then the value of $\Delta P$ is \_\_\_\_\_\_ $Nm^{-2}$.

\sect{Section-II (ii): Numerical Value}

\qn{5} A thin circular coin of mass 10 gm and radius 1 cm is initially in a horizontal xy-plane. The coin is tossed vertically up ($+Z$ direction) by applying an impulse of $\sqrt{\dfrac{\pi}{2}}\times10^{-2}$ N-s at a distance 1/2 cm from its center. The coin spins about its diameter and moves along the $+Z$ direction. By the time the coin reaches back to its initial position, it completes n rotations. The value of n is \_\_\_\_. [Given: The acceleration due to gravity g = 10 $ms^{-2}$]
\qfig{p2_phy_s4_q05_1}

\qn{6} A rectangular conducting loop of length 4 cm and width 2 cm is in the xy-plane, as shown in the figure. It is being moved away from a thin and long conducting wire along the direction $\dfrac{\sqrt3}{2}\hat x+\dfrac12\hat y$ with a constant speed v. The wire is carrying a steady current I = 10 A in the positive x-direction. A current of 10 $\mu A$ flows through the loop when it is at a distance d = 4 cm from the wire. If the resistance of the loop is 0.2 $\Omega$, then the value of v is \_\_\_\_\_\_\_\_\_ $ms^{-1}$. [Given: The permeability of free space $\mu_0=4\pi\times10^{-7}\ NA^{-2}$]
\qfig{p2_phy_s4_q06_1}

\qn{7} A string of length 1 m and mass $4\times10^{-5}$ kg is under tension T. When the string vibrates, two successive harmonics are found to occur at frequencies 750 Hz and 1000 Hz. The value of tension T is \_\_\_\_\_\_\_ Newton.

\qn{8} An incompressible liquid is kept in a container having a weightless piston with a hole. A capillary tube of inner radius 0.1 mm is dipped vertically into the liquid through the airtight piston hole, as shown in the figure. The air in the container is isothermally compressed from its original volume $V_0$ to $\dfrac{50}{51}V_0$ with the movable piston. Considering air as an ideal gas, the height (h) of the liquid column in the capillary above the liquid level in cm is\_\_\_\_\_\_\_. [Given: Surface tension of the liquid is 0.075 $Nm^{-1}$, atmospheric pressure is $10^5\ Nm^{-2}$, acceleration due to gravity (g) is 10 $ms^{-2}$, density of the liquid is $10^3\ kg\,m^{-3}$ and contact angle of capillary surface with the liquid is zero]
\qfig{p2_phy_s4_q08_1}

\qn{9} In a radioactive decay process, the activity is defined as $A=-\dfrac{dN}{dt}$, where N(t) is the number of radioactive nuclei at time t. Two radioactive sources, $S_1$ and $S_2$ have same activity at time t = 0. At a later time, the activities of $S_1$ and $S_2$ are $A_1$ and $A_2$, respectively. When $S_1$ and $S_2$ have just completed their $4^{th}$ and $7^{th}$ half-lives, respectively, the ratio $A_1/A_2$ is \_\_\_\_\_\_\_\_\_.

\qfullbegin
\qn{10} One mole of an ideal gas undergoes two different cyclic processes I and II, as shown in the P-V diagrams below. In cycle I, processes a, b, c and d are isobaric, isothermal, isobaric and isochoric, respectively. In cycle II, processes $a'$, $b'$, $c'$ and $d'$ are isothermal, isochoric, isobaric and isochoric, respectively. The total work done during cycle I is $W_I$ and that during cycle II is $W_{II}$. The ratio $W_I/W_{II}$ is $\dfrac{N}{9}$. The value of N is :- (Take: $\ln2=0.7$)
\qfig{p2_phy_s4_q10_1}
\qfullend

\partsec{2}{CHEMISTRY}
\sect{Section-I (i): Single Correct}

\qn{1} The correct molecular orbital diagram for $O_2$ molecule in the ground state is
\begin{qoptions}
\opt{A}{\qfigi{p2_chem_s1_q01_A}}
\opt{B}{\qfigi{p2_chem_s1_q01_B}}
\opt{C}{\qfigi{p2_chem_s1_q01_C}}
\opt{D}{\qfigi{p2_chem_s1_q01_D}}
\end{qoptions}

\qn{2} Consider the following statements related to colloids.\\
(I) Lyophobic colloids are \textbf{not} formed by simple mixing of dispersed phase and dispersion medium.\\
(II) For emulsions, both the dispersed phase and the dispersion medium are liquids.\\
(III) Micelles are produced by dissolving a surfactant in any solvent at any temperature.\\
(IV) Tyndall effect can be observed from a colloidal solution with dispersed phase having almost same refractive index as that of the dispersion medium.\\
The option with the incorrect set of statements is
\begin{qoptions}
\opt{A}{(I) and (II)}
\opt{B}{(II) and (III)}
\opt{C}{(III) and (IV)}
\opt{D}{(II) and (IV)}
\end{qoptions}

\qn{3} In the following reactions, W, X, Y, and Z are the major products.
\qfig{p2_chem_s1_q03_1}
\qfig{p2_chem_s1_q03_2}
\qfig{p2_chem_s1_q03_3}
\qfig{p2_chem_s1_q03_4}
The correct statement about W, X, Y, and Z is
\begin{qoptions}
\opt{A}{Y has six carbons and it undergoes Cannizzaro reaction.}
\opt{B}{W is a primary alcohol with four carbons.}
\opt{C}{Z is a primary amine with six carbons.}
\opt{D}{X undergoes Kolbe's electrolysis to give a six-carbon product.}
\end{qoptions}

\qn{4} A disaccharide X cannot be oxidised by Fehling solution. The acid hydrolysis of X leads to a laevorotatory solution. The disaccharide X is
\begin{qoptions}
\opt{A}{\qfigi{p2_chem_s1_q04_A}}
\opt{B}{\qfigi{p2_chem_s1_q04_B}}
\opt{C}{\qfigi{p2_chem_s1_q04_C}}
\opt{D}{\qfigi{p2_chem_s1_q04_D}}
\end{qoptions}

\sect{Section-I (ii): One or More Correct}

\qn{5} The complex(es), which can exhibit the type of isomerism shown by $[Pd(NH_3)_2Cl_2]$, is(are)\\
\relax[en = $H_2NCH_2CH_2NH_2$]
\begin{qoptions}
\opt{A}{$[Pt(NH_3)(H_2O)(CN)_2]$}
\opt{B}{$[Pt(NH_3)_2Cl_4]$}
\opt{C}{$[Zn(NH_3)_2Cl_2]$}
\opt{D}{$[Co(en)_2(NO_2)(CN)]^+$}
\end{qoptions}

\qn{6} Atoms of metals x, y, and z form face-centred cubic (fcc) unit cell of edge length $L_x$, body-centred cubic (bcc) unit cell of edge length $L_y$, and simple cubic unit cell of edge length $L_z$, respectively.\\
If $r_z=\dfrac{\sqrt3}{2}r_y$; $r_y=\dfrac{8}{\sqrt3}r_x$; $M_z=\dfrac32M_y$ and $M_z=3M_x$, then the correct statement(s) is (are)\\
\relax[Given: $M_x,M_y$, and $M_z$ are molar masses of metals x, y, and z, respectively.\\
$r_x,r_y$, and $r_z$ are atomic radii of metals x, y, and z, respectively.]
\begin{qoptions}
\opt{A}{Packing efficiency of unit cell of x > Packing efficiency of unit cell of y > Packing efficiency of unit cell of z}
\opt{B}{$L_y>L_z$}
\opt{C}{$L_x>L_y$}
\opt{D}{Density of x > Density of y}
\end{qoptions}

\qn{7} In the following reactions, P, Q, R, and S are the major products.
\qfig{p2_chem_s2_q07_1}
The following three starting materials give products Q, R and S respectively (reading top to bottom):
\qfig{p2_chem_s2_q07_3}
$$\xrightarrow[\text{(ii)}H_3O^\oplus]{\text{(i)}NaOH,H_2O}Q$$
$$\xrightarrow[\text{(ii)}H_2CrO_4]{\text{(i)}H_3O^\oplus,\Delta}R$$
$$\xrightarrow[\text{(iii)Ammoniacal }AgNO_3,H_3O^\oplus]{\text{(i)}Mg,\ dry\ ether\text{ (ii)}CO_2,\text{ then }H_3O^\oplus}S$$
The correct statement(s) about P, Q, R, and S is(are)
\begin{qoptions}
\opt{A}{Compound P and S are same.}
\opt{B}{P, Q, R and S are dicarboxylic acids.}
\opt{C}{Compounds Q and R are the same.}
\opt{D}{R undergo aldol condensation and S does \textbf{not} undergo Cannizzaro reaction.}
\end{qoptions}

\sect{Section-II (i): Paragraph Type}

\paraq{1}{2}
The entropy versus temperature plot for phases $\alpha$ and $\beta$ at 1 bar pressure is given. $S_T$ and $S_0$ are entropies of the phases at temperatures T and 0 K, respectively.
\qfig{p2_chem_s3_para1_2_1}
The transition temperature for $\alpha$ to $\beta$ phase change is 600 K and $C_{P,\beta}-C_{P,\alpha}=1\ Jmol^{-1}K^{-1}$. Assume $(C_{P,\beta}-C_{P,\alpha})$ is independent of temperature in the range of 200 to 700 K. $C_{P,\alpha}$ and $C_{P,\beta}$ are heat capacities of $\alpha$ and $\beta$ phases, respectively.

\qn{1} The value of entropy change, $S_\beta-S_\alpha$ (in $Jmol^{-1}K^{-1}$), at 300 K is \_\_. [Use: $\ln2=0.69$; Given: $S_\beta-S_\alpha=0$ at 0 K]

\qn{2} The value of enthalpy change, $H_\beta-H_\alpha$ (in $Jmol^{-1}$), at 300 K is \_\_.

\paraq{3}{4}
A trinitro compound, 1, 3,5 tris-(4-nitrophenyl) benzene, on complete reaction with an excess of Fe/HCl gives major product, which on treatment with an excess of $NaNO_2$/HCl at $0^\circ C$ provides P as the product. P, upon treatment with excess of $H_2O$ at room temperature, gives the product Q. Bromination of Q in aqueous medium then treatment with $Zn/\Delta$ furnishes the product R. The compound P upon treatment with an excess of aniline under mild acidic conditions gives the product S.

\qn{3} The number of heteroatoms present in one molecule of R is \_\_\_\_\_. [Use: Molar mass (in g $mol^{-1}$): H = 1, C = 12, N = 14, O = 16, Br = 80, Cl = 35.5. Atoms other than C and H are considered as heteroatoms]

\qn{4} The total number of nitrogen atoms present in one molecule of S is \_\_\_\_\_. [Use: Molar mass in g $mol^{-1}$: H = 1, C = 12, N = 14, O = 16, Br = 80, Cl = 35.5. Atoms other than C and H are considered as heteroatoms]

\sect{Section-II (ii): Numerical Value}

\qn{5} $H_2S$ (5 moles) reacts completely with acidified aqueous potassium permanganate solution. In this reaction, the number of moles of water produced is x, and the number of moles of electrons involved is y. The value of (y -- x) is \_\_\_\_.

\qn{6} Among $SO_3^{2-}$, $[SiO_4]^{4-}$, $SO_2Cl_2$, $XeOF_2$, $SOF_4$, $BrF_3$, $[Fe(CO)_4]^{2-}$, $[IO_2F_2]^-$, $[PtCl_3(SnCl_3)]^{2-}$, $XeF_4$, and $HgI_4^{2-}$, the total number of species having $sp^3$ hybridised central atom is \_\_\_\_\_\_\_.

\qn{7} Consider the following molecules: $Br_3O_8$, $F_2O$, $H_2S_4O_6$, $H_2S_5O_6$, and $C_3O_2$.\\
If the sum of the total number of atoms existing in their zero oxidation state in each molecule is X and if the sum of the total number of atoms existing in their positive oxidation state in each molecule is Y then the value of (X+Y) is:-

\qn{8} For $He^+$, a transition takes place from the orbit of radius 105.8 pm to the orbit of radius 26.45 pm. The wavelength (in nm) of the emitted photon during the transition is \_\_\_.\\
\relax[Use: Bohr radius, a = 52.9 pm; Rydberg constant, $R_H=2.2\times10^{-18}$ J; Planck's constant, $h=6.6\times10^{-34}$ J s; Speed of light, $c=3\times10^8\ ms^{-1}$]

\qn{9} 50 mL of 0.2 molal urea solution (density = 1.012 g $mL^{-1}$ at 300 K) is mixed with 250 mL of a solution containing 0.06 g of urea. Both the solutions were prepared in the same solvent. The osmotic pressure (in Torr) of the resulting solution at 300 K is \_\_. [Use: Molar mass of urea = 60 g $mol^{-1}$; gas constant, $R=61$ L Torr $K^{-1}mol^{-1}$; Assume, $\Delta_{mix}H=0$, $\Delta_{mix}V=0$]

\qn{10} The reaction of 4-methyloct-ene (P, 2.52 g) with HBr in the presence of $(C_6H_5CO)_2O_2$ gives two isomeric bromides in a 9:1 ratio, with combined yield of 50\%. Of these, the entire amount of the primary alkyl bromide was reacted with an appropriate amount of diethylamine followed by treatment with eq. $K_2CO_3$ to given a non-ionic product S in 100\% yield. The mass (in mg) of S obtained is \_\_. [Use molar mass (in g $mol^{-1}$): H = 1, C = 12, N = 14, Br = 80]

\partsec{3}{MATHEMATICS}

\sect{Section-I (i): Single Correct}

\qn{1} Let $f:[1,\infty)\to\mathbb R$ be a differentiable function such that $3f(1)=1$ and $3\displaystyle\int_1^x tf(t)\,dt = x^2f(x)-\dfrac{x^3}{3}$, $x\in[1,\infty)$. If $e$ denotes the base of the natural logarithm, then $3f(e^2)$ equals
\begin{qoptions}
\opt{A}{$7e^2$}
\opt{B}{$\ln4+e^2$}
\opt{C}{$4e^2$}
\opt{D}{$e^2-4$}
\end{qoptions}

\qn{2} Consider an experiment of tossing a coin repeatedly until the outcomes of two consecutive tosses are same. Let $\dfrac25$ be probability of a random toss resulting in head. If p denotes the probability that the experiment stops with a tail, then $95p$ is equal to
\begin{qoptions}
\opt{A}{65}
\opt{B}{64}
\opt{C}{63}
\opt{D}{61}
\end{qoptions}

\qn{3} If $\cot^{-1}(y)\in(0,\pi)$ and $2\tan^{-1}y\in(-\pi,\pi)\ \forall\ y\in\mathbb R$, then the integer nearest to the sum of solutions of the equation $\tan^{-1}\dfrac{9}{6y-y^2}+\cot^{-1}\dfrac{6y-y^2}{9}=\dfrac{3\pi}{4}$ for $y>0$ is equal to
\begin{qoptions}
\opt{A}{15}
\opt{B}{6}
\opt{C}{9}
\opt{D}{14}
\end{qoptions}

\qn{4} Let the position vectors of points A, B, C, D are $\vec a=\hat i+2\hat j-5\hat k$, $\vec b=3\hat i+6\hat j+3\hat k$, $5\vec c=17\hat i+16\hat j+35\hat k$ and $\vec d=2\hat i+\hat j+x\hat k$ respectively. Then which of the following options is Correct?
\begin{qoptions}
\opt{A}{If A, B, C, D are coplanar, then $x=-1$}
\opt{B}{If $x=1$, then B, C, D are collinear.}
\opt{C}{If A, B, C, D are coplanar, then BD divides AC in ratio 5 : 4 internally}
\opt{D}{If A, B, C, D are coplanar, then AC divides BD in ratio 1 : 2 internally}
\end{qoptions}

\sect{Section-I (ii): One or More Correct}

\qn{5} Let $M=(a_{ij})$, $i,j\in\{1,2,3\}$ be the $3\times3$ matrix such that $a_{ij}=\begin{cases}2, & \left[\dfrac{j+1}{i}\right]=\dfrac{j+1}{i}\\[2mm] 1, & \left[\dfrac{j+1}{i}\right]\neq\dfrac{j+1}{i}\end{cases}$; $[.]$ denotes the greatest integer function. Then which of the following options is/are Correct?
\begin{qoptions}
\opt{A}{$\text{adj}(M)\begin{pmatrix}1\\2\\3\end{pmatrix}=\begin{pmatrix}1\\2\\3\end{pmatrix}$, where adj(M) represents adjoint of M.}
\opt{B}{There exists a non-zero column matrix $\begin{pmatrix}a_1\\a_2\\a_3\end{pmatrix}\in\mathbb R^3$ such that $M\begin{pmatrix}a_1\\a_2\\a_3\end{pmatrix}=\begin{pmatrix}-a_1\\-a_2\\-a_3\end{pmatrix}$}
\opt{C}{The set $\left\{\begin{pmatrix}x\\y\\z\end{pmatrix}\in\mathbb R^3: M\begin{pmatrix}x\\y\\z\end{pmatrix}=\begin{pmatrix}0\\0\\0\end{pmatrix}\right\}\neq\left\{\begin{pmatrix}0\\0\\0\end{pmatrix}\right\}$.}
\opt{D}{The matrix $\text{adj}(M-5I_3)$ is invertible, where $I_3$ is identity matrix of order 3.}
\end{qoptions}

\qn{6} Let $[.]$ denotes the greatest integer function, $f:(0,1)-\left\{\dfrac13\right\}\to\mathbb R$ be the function defined as $f(x)=[4x]\left(x-\dfrac{1-[x]}{3}\right)^{[2x]}\left(x-\dfrac{[2x]+1}{4}\right)^{3+[x]}$. Then which of the following options is/are Correct?
\begin{qoptions}
\opt{A}{The function f is dis-continuous at exactly two points.}
\opt{B}{There exists exactly two points in $(0,1)$ at which the function f is continuous but \textbf{NOT} differentiable.}
\opt{C}{The function f is \textbf{NOT} differentiable at more than three points in $(0,1)$.}
\opt{D}{The minimum value of f is 0.}
\end{qoptions}

\qn{7} Let S be the set of all twice differentiable functions f from $\mathbb R$ to $\mathbb R$ such that $f''(x)<0\ \forall\ x\in(-2,2)$. For $f\in S$, let $X_f$ be the number of points $x\in(-1,1)$ for which $f(2x)=3x$. Then which of the following options is/are Correct?
\begin{qoptions}
\opt{A}{There exists a function $f\in S$ such that $X_f=0$.}
\opt{B}{For every function $f\in S$, we have $X_f\leq2$.}
\opt{C}{There exists a function $f\in S$, we have $X_f>2$.}
\opt{D}{There does \textbf{NOT} exist any function $f\in S$ such that $X_f=1$.}
\end{qoptions}

\sect{Section-II (i): Paragraph Type}

\paraq{1}{2}
Consider an obtuse angled $\triangle ABC$ in which the difference between the largest and the smallest angle is $\dfrac{\pi}{2}$. Also, the sides of $\triangle ABC$ are in arithmetic progression. The points A, B, C lie on circle $x^2+y^2=1$.

\qn{1} If $\Delta$ denotes the area of $\triangle ABC$, then $(8\Delta)^2$ is equal to

\qn{2} If r denotes the inradius and $\Delta$ denotes the area of $\triangle ABC$, then $\dfrac{(8\Delta)^2}{r}$ is equal to

\paraq{3}{4}
Consider the $8\times8$ square in the figure.
\qfig{p2_math_s3_para3_4_1}
Let $A_1,A_2,\dots,A_{81}$ be the points of intersections of the lines in above figure in some order. We say that $A_i$ and $A_j$ are friends if they are adjacent along a row or along a column. Assume that each of these 81 points have an equal chance of being chosen.

\qn{3} Let $p_i$ denote the probability that a randomly chosen point has `$i$' friends, $i=0,1,2,3,4$. Let X denote the random variable such that for $i=0,1,2,3,4$, the probability that $P(X=i)=p_i$. Then value of $36E(X)$ is equal to

\qn{4} Two distinct points are chosen from the given 81 points. Let p denote the probability that they are friends. Then the value of $36p$ is equal to

\sect{Section-II (ii): Numerical Value}

\qn{5} For $x\in\mathbb R$, let $|2\tan^{-1}x|<\pi$. Then the minimum value of the function $f:\mathbb R\to\mathbb R$ defined by $f(x)=\displaystyle\int_0^{x\tan^{-1}x}\left(\dfrac{e^{t-\cos t}}{1+t+t^2+\dots+t^{2023}}\right)dt$

\qn{6} For $x\in\mathbb R$, let $y(x)$ be a solution of the differential equation $(x^4-5)\dfrac{dy}{dx}=4x^3y-4x^3(x^4-5)^2$, $y(\sqrt2)=7$. Then the maximum value of the function $y(x)$ is

\qn{7} Let X be the set of all five digit numbers formed using 1, 1, 2, 2, 2, 4, 4, 0. For example, $22240\in X$, but $02244\notin X$ and $44422\notin X$. Suppose that each element of X has an equal chance of being chosen. Let p denote the conditional probability that an element of X chosen at random is a multiple of 20, given that it is a multiple of 5. Then the value of $310p$ is equal to

\qn{8} Let $A_1,A_2,A_3,\dots,A_{12}$ be the vertices of a regular dodecagon that lie on a circle with center at $(0,0)$ and radius 5 units. Let P be a point on the circle and $PA_i$ denote the distance between the points P and $A_i$ for $i=1,2,\dots,12$. Then the maximum possible value of the product $\left(\left(\displaystyle\prod_{r=1}^{12}(PA_r)\right)^2\Big/2^{23}\right)^{\frac{1}{12}}$ is equal to

\qn{9} Consider the following set $A=\left\{\begin{pmatrix}a&3&b\\c&2&d\\0&5&0\end{pmatrix}:a,b,c,d\in\{0,1,2,3,5,7,11\}\right\}$. If $\left\{\begin{pmatrix}x\\y\\z\end{pmatrix}\in\mathbb R^3: M\begin{pmatrix}x\\y\\z\end{pmatrix}=\begin{pmatrix}0\\0\\0\end{pmatrix},M\in A\right\}=\left\{\begin{pmatrix}0\\0\\0\end{pmatrix}\right\}$, then the number of matrices in the set A is

\qn{10} Let $S_1$ be a circle of radius 1 and center $(-1,3)$. Let $S_2$ be a circle of radius r and with center at $A=(5,-7)$. Points $P_1,Q_1$ lie on $S_1$ and $P_2,Q_2$ on $S_2$ such that $P_1P_2$ is tangent to both $S_1$ and $S_2$. Also $Q_1Q_2$ is tangent to both $S_1$ and $S_2$. The line joining mid-points of $P_1P_2$ and $Q_1Q_2$ meets X-axis at B. If $AB=\dfrac{\sqrt{53}}{2}$, then the value of $r^2$ is

\end{multicols}
\end{document}

